\section{模型检验与讨论}
\subsection{检验体系归纳}
本文的检验不是在得到结果后选择有利指标，而是针对不同模型结构预先采用五类检查：解析极限用于检查公式是否退化正确；结构对称性用于检查方向模型；边界残差用于检查递推与覆盖条件；独立复算用于检查路线数据和评估器一致性；网格敏感性用于判断离散覆盖结论是否稳定。表~\ref{tab:validation-summary} 汇总检验对象和结论。

\begin{table}[H]
\centering
\caption{各问题模型检验汇总}
\label{tab:validation-summary}
\begin{tabularx}{\textwidth}{C{1.1cm}C{3.4cm}L L}
\toprule
问题 & 检验方法 & 具体检查量 & 结论 \\
\midrule
一 & 零坡度极限、单调性 & $\alpha=0$ 时 $W=2D\tan\varphi$；$W$ 随 $D$ 单调增加 & 解析退化正确，表格趋势一致 \\
二 & 对称性与不变量 & $45^\circ/315^\circ$、$135^\circ/225^\circ$ 对称；$90^\circ/270^\circ$ 沿程常值 & 全部满足至浮点精度 \\
三 & 边界和递推残差 & 首线西边界残差、$\eta_i-0.10$、第 33/34 条东边界差 & 34 条可行，33 条在路线族内不可行 \\
四 & 几何、独立复算、敏感性 & 长度差、越界、交叉、硬拐角、曲率半径、三组网格漏测率 & 几何通过；高分辨率网格暴露微小漏测 \\
\bottomrule
\end{tabularx}
\end{table}

\subsection{问题一和问题二的一致性检验}
问题一在 $\alpha=0$ 时有 $L_d=L_s=D\tan\varphi$，说明斜坡公式与常用平底覆盖公式衔接。对于给定 $\alpha$，式~\eqref{eq:p1-half} 中的比例系数保持常数，因此覆盖宽度必须与水深成正比；表~\ref{tab:p1-result} 的相邻宽度差保持近似常量，与解析结构相符。问题二中，$D(l,\beta)$ 对 $\cos\beta$ 敏感，而 $\gamma(\beta)$ 对 $|\sin\beta|$ 敏感。将 $\beta$ 替换为 $360^\circ-\beta$ 时这两个量不变，因此对应结果必须相等；数值矩阵的四组对称关系均严格成立。上述检验能发现角度方向、海里换算或正负号错误，因而比只比较单个数值更有效。

\subsection{问题三的边界残差与最少性复核}
对 $34$ 条测线逐条代回式~\eqref{eq:p1-overlap}，最小和最大相邻重叠率均为 $0.100000000000$；第一条线深水边界与西边界的误差小于 $10^{-9}$ m。第 $33$ 条线浅水边界距东边界尚有 \,\SI{25.756620}{m}，第 $34$ 条则超过 \,\SI{15.405496}{m}。由于显式递推每一步均选取允许的最大间距，这两个边界量共同构成路线族内最少性证据，而不只是“程序输出 34”这一孤立结果。

\subsection{问题四的几何与独立复算}
显式几何模块按直线精确长度和 Bézier 段 Simpson 积分求和，得到 \,\SI{576726.635512}{m}；独立空间评估器使用覆盖采样折线求和，得到 \,\SI{576725.729283}{m}。两者差值 \,\SI{0.906230}{m} 仅占总长度的 $1.57\times10^{-6}$，说明曲线采样分辨率足以支持空间指标计算。路线越界长度为 $0$，硬拐角数为 $0$，相邻线交叉数为 $0$，并且所有 Bézier 段导数均非零。最小曲率半径 \,\SI{225.983}{m} 是几何诊断值，不与未给出的船舶阈值比较。

\subsection{网格敏感性与不利结果讨论}
推荐方案分别在 $30\times30$、$50\times50$ 和 $70\times70$ 方格中心上复算。前两组漏测比例为 $0$，$70\times70$ 网格出现 $0.040816\%$ 漏测；多重覆盖比例分别为 $48.777778\%$、$48.000000\%$ 和 $48.387755\%$。图~\ref{fig:p4-sensitivity} 表明，漏测结论对网格分辨率存在轻微敏感性，而多重覆盖比例在约 $48\%$ 附近相对稳定。该微小漏测很可能位于条带边缘或曲线局部弯折附近，它没有显著改变总体覆盖水平，但足以否定“连续区域严格无漏测”的表述。

\begin{figure}[H]
\centering
\includegraphics[width=0.68\textwidth]{problem4_sensitivity.pdf}
\caption{推荐方案的漏测率网格敏感性；$70\times70$ 网格暴露微小漏测}
\label{fig:p4-sensitivity}
\end{figure}

因此，本文把最终结论限定为“名义冻结网格漏测率为 $0$，高分辨率敏感性网格出现 $0.040816\%$ 漏测”。若实际工程要求任何采样尺度都不得出现漏测，应采用 A 方案或在 C-05 的局部缺口附近增加补测线，而不能忽略敏感性结果。
